\documentclass[11pt,a4paper]{article}
\usepackage{bookmark}
\usepackage[utf8]{inputenc}
\usepackage[english]{babel}
\usepackage{graphicx}
% \usepackage{multicol}


\begin{document}
\begin{multicols}{2}
\author{Gonz\'alez Tapia Jaime}
\title{M\'odulo de Young}
\date {13 de septiembre 2015}
\maketitle

\section{Introducci\'on}
El m\'odulo de Young  es un par\'ametro que caracteriza el comportamiento de un material el\'astico, seg\'un la direcci\'on en la que se aplica una fuerza. Para un material el\'astico lineal e is\'otropo, el m\'odulo de Young tiene el mismo valor para una tracci\'on que para una compresi\'on, siendo una constante  independiente del esfuerzo siempre que no exceda de un valor m\'aximo denominado l\'imite el\'astico, y es siempre mayor que cero: si se tracciona una barra, aumenta de longitud, no disminuye. Este comportamiento fue observado y estudiado por el científico ingl\'es Thomas Young.


Tanto el m\'odulo de Young como el l\'imite el\'astico son distintos para los diversos materiales. El m\'odulo de elasticidad es una constante el\'astica que, al igual que el l\'imite el\'astico, puede encontrarse emp\'iricamente con base al ensayo de tracci\'on del material.


\section{Trabajo de clase}
El trabajo en clase consist\'ia en obtener la compresi\'on y la cantidad de cent\'imetros o metros cuadrados que se iba a reducir el hueso, tomando como referencia el trabajo de clase se plantea el hacer lo mismo, pero sustituyendo la masa con la propia del individuo, al ser mi masa del total de 725 Newtons, la gr\'afica queda de la siguiente forma:
\end{multicols}

\begin{multicols} {2}


\section{Conlclusiones}
El ejemplo de teo\'ia y la tarea era en mi caso casi id\'entico, solamente variaba por 75, que son b\'asicamente nada, pero descubr\'i que estoy en pedo "adecuado" para no doblegar mis huesos al punto de quiebre.
\end{multicols}
 
\section{Bibliograf\'ia}
\url{http://fisica-albarracin.es.tl/M%F3dulo-de-Young-.htm}

\url{https://es.sharelatex.com/learn/Multiple_columns#Column_separation}

\url{https://www.overleaf.com/help/111-how-do-i-insert-an-image-into-my-document#.VfY1CZYzLYU}


\end{document}
\author{Gonzalez Tapia Jaime}
\title{M\'odulo de Young}
\date {13 de septiembre 2015}